% 全文基本假设与符号说明。
% 本节放在摘要之后、问题一正文之前，用于统一四问的变量、数据接口和适用边界。
\section{基本假设与符号说明}
\label{sec:assumptions}

本文研究对象横跨文本质量评价、领域配比、训练损失、模型规模和算力预算。不同问题使用的数据块并非同一批训练实验，因此本节先统一符号，再明确哪些关系来自观测数据、哪些关系只用于条件情景分析。下列假设是模型的工作前提；它们不等同于已经由实验验证的结论。

\subsection{基本假设}

\begin{enumerate}
  \item \textbf{数据角色与证据分级假设。} A1--A3 用于多源文本质量评价，A4--A15 用于领域配比与交叉熵损失分析，B1--B10 用于规模、质量和跨来源迁移分析，C 数据用于问题四的评测前沿分析。半合成、估算或外推数据只用于相应的条件分析和敏感性审计，不被当作与真实训练实验等价的独立观测。

  \item \textbf{质量指标方向假设。} 原始质量字段经过字段级转换后形成 16 项可入模指标。句末标点率、整洁度、无广告程度、可读性、流畅度、独特词比例、教育性、写作、事实性及目标域亲和度等按正向指标处理；无字母词比例和最高频二元、三元词组占比按负向指标处理，作为噪声或重复性的代理。该方向约定可能对代码、数学、表格、LaTeX 和非英语文本产生适用性偏差，因此只表示建模选择。

  \item \textbf{冻结标尺假设。} A1 拟合集上的 $1\%$ 与 $99\%$ 分位数用于稳健归一化，CRITIC 权重也只在该拟合集上估计；留出集、A2 和 A3 不重新拟合这些参数。这样可以保证不同数据分区在同一评分标尺下比较，但不表示扩展集与 A1 具有完全相同的分布。

  \item \textbf{质量聚合与冲突假设。} 样本级质量分数先计算、后按记录均值聚合为领域级和语料级分数。16 项指标的 TOPSIS 总分采用完全可补偿聚合；冲突分析将通用指标压缩为内容价值、表达质量、整洁度、无广告程度和非重复程度五维，并以 A1 拟合集的 P20/P80 相对阈值定义候选冲突。有限补偿参数 $\lambda\in[0,1]$ 表示对最弱维度的保守程度，不被解释为由人工真值估计出的最优参数。

  \item \textbf{单纯形与配比假设。} 17 个领域配比非负且和为 1，即属于单纯形 $\Delta_{17}$。为避免闭合约束造成的虚假相关，配比模型在带零替换的 ILR 坐标中拟合；零替换常数和 ILR 基会影响数值系数，因此配比系数只在声明的预处理协议下解释。

  \item \textbf{尺度分层假设。} 同一参数规模内可以研究配比与损失的响应关系，但只含配比的模型不能同时解释同一配方在不同规模下的绝对损失。因而跨规模比较必须显式保留规模信息；A6/A8、A12/A14 等配对结果主要用于分离规模效应与配方迁移，不能据此建立无条件的统一规模律。

  \item \textbf{跨来源接口假设。} Q1 的领域质量代理与 B6/B7 的原生质量分数处于不同量尺，A、B 数据也缺少行级联合连接键。问题二、问题三中的 $Q_A\mapsto Q_B$ 映射、迁移强度和配比残差化只构成显式的条件接口；它们不被解释为已经识别的质量因果效应、跨来源标定关系或四变量联合定律。

  \item \textbf{损失与预算假设。} 交叉熵验证损失 $L$ 是 A、B 数据块之间的共同响应变量。问题三中，参数规模、训练数据量和上下文长度的基础计算成本按给定的乘积形式计算，质量处理成本按预设的指数、幂函数或对数成本族进行敏感性分析。上下文长度由架构与任务条件外生给定，不作为连续内点变量；支持域之外的优化解只作外推诊断。

  \item \textbf{问题四预测假设。} 问题四以可比时间窗口内的 C8 六任务等权分数和月度横截面高分位数为预测对象。滚动基线和面板分位数模型用于时间外预测评价；问题三的规模弹性只作为算力情景的外生输入。时间项可能同时吸收技术进步、模型构成和评测规则变化，因此不直接解释为纯技术进步。
\end{enumerate}

\subsection{符号说明}

表~\ref{tab:global-notation} 汇总全文反复使用的变量、集合、指标和接口。下标 $i$、$r$、$j$、$g$、$t$ 分别表示样本、配方行、质量或损失指标、领域/质量维度和时间（月）；未特别说明时，向量采用列向量记号，$\boldsymbol 1$ 表示全 1 向量。

\begin{table}[htbp]
  \centering
  \caption{全文主要符号、含义及使用范围}
  \label{tab:global-notation}
  \small
  \renewcommand{\arraystretch}{1.18}
  \begin{tabular}{p{0.19\linewidth}p{0.52\linewidth}p{0.19\linewidth}}
    \toprule
    符号 & 含义 & 主要使用部分 \\
    \midrule
    $A_1,A_2,A_3$ & 文本质量抽样集与扩展集；合并去重后形成唯一记录集合 & 问题一 \\
    $A_4$--$A_{15}$ & 17 域配方及 13 个交叉熵损失的拟合、验证和外推表 & 问题一第三小问 \\
    $B_1$--$B_{10}$ & 模型规模、训练量、质量扰动和迁移审计数据块 & 问题二、问题三 \\
    $C_1$--$C_8$ & 模型评测、时间序列、架构元数据和逐任务评测数据块 & 问题四 \\
    $x_{ij}$ & 第 $i$ 条文本在第 $j$ 个原始质量字段上的取值 & 问题一 \\
    $\widetilde{x}_{ij}$ & 经过正负方向统一后的质量指标 & 问题一 \\
    $z_{ij}$ & 基于 A1 拟合集分位边界归一化后的质量指标，取值在 $[0,1]$ & 问题一 \\
    $w_j$ & CRITIC 方法得到的第 $j$ 个指标权重，满足 $\sum_jw_j=1$ & 问题一 \\
    $Q_i$ & 第 $i$ 条文本的 TOPSIS 样本级质量分数，范围为 $[0,100]$ & 问题一 \\
    $Q_g$、$Q_{\mathrm{corpus}}$ & 领域级质量均值与全语料质量均值 & 问题一及后续接口 \\
    $G_g,z_{ig},v_g$ & 五个质量维度的指标集合、组内加权分数和组权重 & 冲突消解 \\
    $H_i,L_i,I_i$ & 第 $i$ 条文本的高分维度集合、低分维度集合和候选冲突标记 & 冲突消解 \\
    $C_i,B_i$ & 冲突强度与高低维度配对广度 & 冲突消解 \\
    $M_i,m_i,\lambda$ & 五维加权均值、最弱维度分数和有限补偿参数 & 冲突消解 \\
    $\boldsymbol p$、$\Delta_{17}$ & 17 域配比向量及其单纯形约束，$p_j\ge0,\sum_jp_j=1$ & 配比建模 \\
    $\boldsymbol z_r$ & 第 $r$ 行配比经 ILR 变换后的坐标 & 配比建模 \\
    $y_{rj}$、$\widehat y_{rj}$ & 第 $r$ 行第 $j$ 个损失及其模型预测值 & 配比建模 \\
    $N,D$ & 参数量和训练数据量，均以十亿为单位 & 问题二、问题三 \\
    $L$ & 验证集交叉熵损失；$L_j$ 表示第 $j$ 个损失目标 & 问题一至问题三 \\
    $E,A,B,\alpha,\beta$ & 幂律规模项的截距、系数和指数 & 问题二、问题三 \\
    $Q_A,Q_B,T(\cdot)$ & A 侧配方质量代理、B 侧原生质量坐标及两者之间的假设映射 & 问题二、问题三 \\
    $\boldsymbol p_0,\boldsymbol v,\boldsymbol r(\boldsymbol p)$ & 参考配方、质量方向载荷和配比残差 & 问题二 \\
    $\widehat{\boldsymbol w},\lambda_Q,\lambda_p$ & 配比响应载荷与质量/配比跨来源迁移强度 & 问题二、问题三 \\
    $C,L_{\mathrm{ctx}}$ & 总算力预算（FLOPs）与上下文长度 & 问题三、问题四情景 \\
    $\eta,g_k(Q)$ & 注意力成本系数与第 $k$ 类质量处理成本函数 & 问题三 \\
    $S_i,F_t$ & 模型 $i$ 的 C8 六任务综合分数与月份 $t$ 的前沿分数 & 问题四 \\
    $N_i,t_i,\tau$ & 模型参数规模、相对月份和条件分位数水平 & 问题四 \\
    $\mathrm{RMSE},\mathrm{MAE}$ & 均方根误差和平均绝对误差 & 模型评价 \\
    \bottomrule
  \end{tabular}
\end{table}

数据块标签 $B_1$--$B_{10}$ 与规模律中的系数 $B$、冲突广度 $B_i$ 仅在不同上下文中使用；阅读公式时，以下标或粗体出现的量优先按变量定义解释。

为避免不同问题之间的量尺混用，本文特别约定：$Q_i$ 和 $Q_g$ 使用 0--100 的质量分数；B6/B7 的原生 $Q_B$ 使用其自身的 $[0.1,1.0]$ 坐标；$L$ 为交叉熵损失，数值越低表示模型在对应验证口径下的损失越小；$p$ 是配比向量而不是质量分数。除非另有说明，问题二和问题三中的联合表达式均指冻结参数下的条件预测，不表示新增的联合训练观测。
